\section{Measurement Model}

Let \(\boldsymbol{\phi}(t)\) denote a vector of locally extracted facial-behavior features at time \(t\). Candidate components include landmark geometry, head pose, eye-region visibility, selected blendshape dynamics, and tracking-quality measures. Identity embeddings and face-recognition features are excluded.

For participant \(i\), each feature is standardized relative to an individual baseline:

\[
z_{ij}(t)=
\frac{\phi_{ij}(t)-\mu^{base}_{ij}}
{\sigma^{base}_{ij}+\epsilon}.
\]

A descriptive deviation index may be defined as

\[
D_{\phi,i}(t)
=
\sqrt{\sum_j w_j z_{ij}(t)^2},
\qquad w_j \ge 0.
\]

This index is not interpreted as valence, arousal, anxiety, or any categorical emotion. It merely quantifies distance from a participant-specific behavioral reference pattern.

The camera channel is evaluated against a base information set

\[
\mathbf{b}_t =
[
\mathbf{P}_t,
\mathbf{I}_t,
\mathbf{S}^{self}_t,
\mathbf{C}_t
],
\]

where task performance, optional inertial summaries, explicit child input, and task context are known. The extended information set is

\[
\mathbf{b}^{+}_t = [\mathbf{b}_t,\boldsymbol{\phi}(t)].
\]

The scientific contribution of the camera exists only if prediction or decision quality improves when moving from \(\mathbf{b}_t\) to \(\mathbf{b}^{+}_t\).
